% sec-05: the evidence a funder is buying.  Table 10.
\section{The Evidence Behind the Robotic Phase 1}

The request to the SEC office concerns the company's research as a body of
work. A funder reading this packet has a narrower question: what trial would the
money run, and what evidence says it is worth running? This section answers with
six quantities a reviewer can check, each beside the limitation its own authors
stated, and with the design numbers fixed in the deposited protocol.

\begin{apptable}
\begin{tabularx}{\textwidth}{>{\raggedright\arraybackslash}p{3.6cm} >{\raggedright\arraybackslash}p{2.8cm} >{\raggedright\arraybackslash}p{2.7cm} >{\raggedright\arraybackslash}p{1.7cm} >{\raggedright\arraybackslash}X}
\headrow \thead{Source} & \thead{Test arm} & \thead{Comparator} & \thead{Tier} & \thead{Stated limitation}\\
\midrule
RASolute 302, 2026 & 13.2 mo mOS & 6.6 mo mOS & Trial & Metastatic setting only \\
QSP model, 250 ODEs & 12.8 mo, HR 0.25 & 5.4 mo mOS & In silico & No resistance modeled \\
Twin, 1000 patients & 12.1 mo mOS & No comparator & Twin & Simple Emax model \\
Twin, PFS & HR 0.31 & No comparator & Twin & No immune compartments \\
VVUQ, 55 tests & Score 81.9 & V and V 40 gate & Twin & Pre-trial score only \\
Empirical triplicate & G3+ 8.0 pct & G3+ 25.0 pct & In silico & Two logs disagree \\
\end{tabularx}
\end{apptable}
\vspace{-0.60cm}
\tabcap{Table 10. Six quantities that a reviewer can check against a published or deposited\\
source, each set beside its comparator, its evidence tier, and its stated limitation.}

Only the first row is a clinical result, and it is the developer's, in the
metastatic setting \cite{rasolute302}. The other five are the company's own
in silico and digital twin work \cite{qspmetpancre,daraxsims}. No row is
presented without its limitation, and none is presented as evidence for the
perioperative use this program proposes. What the six do establish is that the
agent choice was made on a dated, quantitative record, and that the one
clinical readout since has landed near the company's simulation.

\subsection{The trial the evidence supports}

The deposited Phase 1 protocol is open-label, single-arm, 3+3 dose escalation of
perioperative daraxonrasib at 160, 220 and 300 mg, with a 28-day dose-limiting
toxicity window, in up to 18 treated participants with resectable or
borderline-resectable KRAS-mutated pancreatic ductal adenocarcinoma
\cite{onpremwhippl}. The procedure is a staged eight-arm robotic
pancreaticoduodenectomy with 56 degrees of freedom and 640 sensor channels, force
limits of 3 N per arm and 18 N cumulative, a 3 ms cross-arm stop against a 500 ms
system stop, and a five-vessel no-fly gate. An on-premises language model is
confined to advisory output, with no write credential to data capture and no
route to the robot control network \cite{phase1ind}.

The disease the trial addresses remains the one with the widest gap between
incidence and survival in the United States \cite{siegel2025statistics}. The
developer's own adjuvant Phase 3 in resected disease is enrolling
\cite{ctgov304}. The perioperative, robot-assisted question this program asks is
a different one from either of the developer's Phase 3 trials, and it is the
question a Phase 1 is designed to open.

\subsection{The money frame}

The program's direct cost is \$700,000 a year for five years, \$3,500,000 in
total. An SBIR route of \$306,000 in Phase I and \$1,300,000 in Phase II totals
\$1,606,000, leaving a delta of \$2,104,000 that the capitalization plan places
against \$5,900,000 of private capital behind a firewall, a leverage of 3.67 to 1
\cite{capitalplan2026,tenapps2026}. The company's comparable virtual trial work is
\textbf{projected} at \$36,330 per run, a figure that describes simulation work
already deposited and not the cost of the clinical trial. None of these numbers
depends on the SEC request, and none changes with its outcome.
